\chapter{Abstract}
\begin{tabular}{@{} >{\bfseries}p{0.2\textwidth} p{0.77\textwidth} @{}}
\MakeUppercase{Keywords} & Enhanced Voltage Protection Module ($EVPM$), Hierarchical Microgrid Architectures, Transform-Domain Deep Convolutional Neural Networks ($CNN$)\\ \end{tabular}
\par










\begin{doublespacing}

The global electrical infrastructure is undergoing a fundamental structural transition from centralized, synchronous-machine-dominated architectures to decentralized, power-electronic-interfaced networks. This research asserts that such a shift renders legacy deterministic protection paradigms—specifically fundamental-frequency, RMS-based overcurrent (ANSI 51), directional (ANSI 67), and distance (ANSI 21) elements—fundamentally obsolete. These conventional methodologies exhibit existential susceptibility to protective blinding and catastrophic under-reach when subjected to the non-linear, converter-limited transient signatures inherent to modern networks. Furthermore, existing frameworks fail to resolve critical operational paradoxes introduced by modern grid codes, notably the Low Voltage Ride Through (LVRT) and dynamic reactive current injection mandates of IEEE 1547-2018 and IEEE 2800. To definitively overcome these limitations, this thesis proposes a comprehensive, multi-domain protective architecture that synthesizes

\begin{itemize}
\item \textbf{ Algorithms for Multi -Functional Adaptive Interconnection Protection for IBRs:} Detail the formulation and validation of instantaneous time-domain sampling logics for frequency-independent fault detection.

\item \textbf{ Adaptive protection architectures for microgrids:} Presents the multi-layered communication-assisted coordination strategy for grid-tied and islanded transitions.

\item \textbf{AI/ML-based advanced protection for future microgrid:} Introduces the PINN-based (Physics-Informed Neural Network) architecture for robust fault detection in low-inertia environments. 

\end{itemize}
Through rigorous analysis of the simulation data, the performance of the proposed algorithms is compared with that of deterministic legacy relays across a spectrum of extreme operational scenarios.


\end{doublespacing}
